\section{Related Work}
\label{sec:related}

\subsection{Cache Management for LLM Inference}
\label{sec:related:cache}

Production inference engines optimize transformer serving through paged memory allocation. Systems like PagedAttention and vLLM \cite{kwon2023pagedattention} manage the KV cache by dividing sequences into fixed-size blocks, reducing fragmentation and enabling efficient scatter-gather memory access. High-performance kernel libraries, such as FlashInfer, and production execution engines, like TensorRT-LLM, adopt similar block-based abstractions. These systems handle a single, homogeneous cache type exceptionally well. They assume uniform memory access patterns and predictable block scaling tied directly to sequence length. AVMP extends this page-based paradigm to support two physically heterogeneous cache types simultaneously. Rather than replacing these highly optimized attention managers, AVMP provides a complementary memory abstraction that integrates with existing page tables to support hybrid model execution.

\subsection{Hybrid Architecture Serving}
\label{sec:related:hybrid}

Recent systems extend inference engines to support hybrid architectures containing both State Space Model and transformer layers \cite{lieber2024jamba, dao2024mamba2}. SGLang provides production serving for these models using a static dual-pool approach \cite{zheng2024sglang}. It provisions a \texttt{HybridReqToTokenPool} for attention and a \texttt{HybridLinearKVPool} for SSM state. Operators configure the partition at engine startup via the \texttt{mamba\_full\_memory\_ratio} parameter. This rigidity causes capacity failures when prompt distributions shift. As we demonstrate in \S\ref{sec:eval}, a default 0.9 ratio triggers 1221 OOM events compared to 552 OOMs at a 0.5 ratio on identical workloads.

Alternatively, vTensor introduces GPU virtual memory management for KV caches using hardware features to decouple physical memory mapping from virtual address spaces \cite{xu2024vtensor}. While vTensor targets a single cache type to reduce fragmentation, AVMP applies this virtual abstraction to two heterogeneous pools. AVMP contributes a unified virtual handle system combined with runtime capacity rebalancing, addressing both the correctness requirements and the adaptivity challenges in hybrid serving.

\subsection{Dynamic Resource Allocation}
\label{sec:related:dynamic}

Inference engines employ runtime adaptation to maximize hardware utilization. Continuous batching strategies adapt batch sizes dynamically at iteration boundaries to increase throughput. AVMP adapts pool capacity at request-level boundaries, using a more conservative trigger placed strictly in the allocation path. Memory eviction strategies, such as sliding window attention, discard older tokens dynamically under memory pressure. These techniques operate within a single pool; AVMP migrates capacity across two heterogeneous pools.

At the framework level, tools like the PyTorch \texttt{expandable\_segments} allocator manage dynamic segment growth directly in CUDA. This allocator operates below the inference engine abstractions. AVMP operates at the pool level, translating application-level \texttt{CapacityError} exceptions into physical capacity migrations between distinct cache implementations. AVMP's CapacityError-triggered migration is conservative: it fires only when allocation would otherwise fail, not preemptively on transient pressure.
